\documentclass[12pt]{article}
\usepackage[utf8]{inputenc}
\usepackage[T1]{fontenc}
\usepackage[margin=1in]{geometry}
\usepackage{mathpazo}
\usepackage{setspace}
\usepackage{parskip}
\usepackage{hyperref}
\hypersetup{hidelinks}

\onehalfspacing

\title{\textbf{Case Reflection 2: Duolingo}\\[0.3em]
\large Duolingo: Language Learning Case Reflection}
\author{Christian Nikolov}
\date{March 22, 2026}

\begin{document}
\maketitle
\thispagestyle{empty}

Duolingo is popular because it costs nothing, runs on your phone and each lesson takes three to five minutes. Von Ahn gave a TED Talk that got watched by two million people and that alone seeded the user base. From there until 2019 they didn't do any advertising or marketing at all. Apple named it iPhone App of the Year in 2013. By 2021 they had 40.5 million monthly active users and over half a billion downloads across 40+ languages, including Klingon and Dothraki. In the U.S. alone more people were learning on Duolingo than in the entire public education system. Von Ahn figured out that if you make language learning feel like a game instead of homework, people will actually do it.

The two products could not be more different. Rosetta Stone started in 1992 selling CD-ROM software in large yellow boxes and still feels like it. By 2022 they'd moved to online subscriptions at \$8 to \$12 a month but the format is the same: 30-minute structured lessons, serious tone, no games. Duolingo is free, sessions last three to five minutes and there's an animated owl that guilts you into coming back. Rosetta Stone had about half a million subscribers and \$200 million in annual revenue. Duolingo had 500 million downloads. Rosetta Stone is trying to teach you properly. Duolingo is trying to get you hooked. A private equity firm bought Rosetta Stone for roughly \$800 million in 2021, the same year Duolingo went public at a much higher valuation.

For customer acquisition, von Ahn said that from launch until 2019 ``our growth was solely due to positive word of mouth, we didn't do any advertising or marketing at all.'' The Wall Street Journal and Forbes called it the best free language app, which helped. After 2019 they started spending on performance marketing in Japan, India and Southeast Asia. Sales and marketing went from \$14.9 million in 2019 to \$59.1 million in 2021. But even with real ad spend, most growth was still organic.

I don't think they're really competing for the same people. Duolingo is going after the estimated two billion people worldwide who want to learn a language but can't or won't pay. Von Ahn said it himself: ``Knowledge of English in a non-English speaking country can usually mean that your income potential is doubled.'' He built the app for someone in Guatemala, not someone in Connecticut. Rosetta Stone's audience is willing to spend \$8 to \$12 a month on structured instruction. Different price, different motivation, different product.

Retention is where Duolingo is strongest and also where I have the most questions. Streaks work. People post thousand-day streaks on social media and protect them like they're real achievements. Leaderboards ``single-handedly increased the time people spent on Duolingo by 20\%'' according to von Ahn. Hearts punish mistakes so you actually focus. The notifications are so aggressive that users make memes about them. Von Ahn admitted they can be ``perceived as overly-aggressive'' but justified it: ``If you want to learn a language, you have to turn it into a habit.'' All of that keeps people coming back. What I'm less sure about is whether coming back means learning. One New York Times reviewer put it bluntly: ``Language apps are not other humans.'' You can tap through exercises for years and never hold a conversation. The company itself says finishing a course gets you to ``advanced beginner or early intermediate.'' That's honest but it's also a ceiling. And with 94\% of users on the free tier, Duolingo needs the gamification to eventually push people toward paying, not just playing.

\vfill
\noindent\small\textit{Reference: Moon, Y. (2022). Duolingo: Teaching Languages to the Masses. Harvard Business School Case 9-323-016.}

\vspace{1em}
\noindent\footnotesize\textit{Word count: 600 | Compiled: \today}

\end{document}
